% Run GCT_KRON_20261120; GPT-6 Astra Pro; effort not exposed.
\section{A positive rule for three equal height-three rectangles}

\subsection*{A-026. An actual Kronecker polynomial ring with a saturation hole}
For every integer $h\geq0$, with the empty-partition convention
at $h=0$, one has the positive formula
\[
 \boxed{\quad
 g((h,h,h),(h,h,h),(h,h,h))
    =\#\{(a,b,c)\in\mathbb Z_{\geq0}^3:2a+3b+4c=h\}.
 \quad}                                                 \tag{A-026.1}
\]
In particular the coefficient can be evaluated by the closed form
\[
 \begin{cases}
 \left\lfloor\dfrac{h^2+12h+48}{48}\right\rfloor,&h\text{ even},\\[3pt]
 \left\lfloor\dfrac{h^2+6h+21}{48}\right\rfloor,&h\text{ odd}.
 \end{cases}                                             \tag{A-026.2}
\]
This is an extraction from a known invariant-ring theorem,
not a claim of a new general Kronecker rule. It supplies
positive monomial bases in an unbounded family, with
unbounded multiplicities, and gives a precise limit to the
multiplicity-free construction in A-024.

\paragraph{Exact primary input and hypotheses.}
Let $E,F,G$ be complex vector spaces of dimension three,
and let
\[
       \mathcal A=\mathbb C[E\otimes F\otimes G]
                  ^{\mathrm{SL}(E)\times\mathrm{SL}(F)\times\mathrm{SL}(G)}.
\]
The established invariant-theory result is
\[
                  \mathcal A=\mathbb C[I_6,I_9,I_{12}],           \tag{A-026.3}
\]
with algebraically independent generators of the displayed
ordinary polynomial degrees. A precise primary locator is
Domokos--Drensky, \emph{Defining relation for semi-invariants
of three by three matrix triples}, arXiv:1101.3178v1
(17 January 2011), Section 2, Theorem 1 and Remark 2(ii),
printed p.~6. The characteristic-zero hypothesis is stated
at the start of Section 2, p.~3. Its generators are
$I_6=H$ from (3), p.~4, $I_9=Q$ from (4), p.~5,
and $I_{12}=\widetilde S$, defined by the ternary-cubic
substitution on p.~5. All these pages were read in the
primary PDF; pp.~5--6 were also visually checked.

The algebraic step in that proof is transparent. Its invariant
algebra is free with basis $1,Q$ over
$\mathbb C[\widetilde S,\widetilde T,H]$, and Theorem 1 gives
\[
                 Q^2=H^3+27H\widetilde S-\frac{27}{4}\widetilde T.
\]
The free-module statement makes this the only relation.
Since the coefficient of $\widetilde T$ is nonzero in
characteristic zero, elimination of $\widetilde T$ gives
the polynomial algebra in $H,Q,\widetilde S$.
This checks the deduction used in (A-026.3). The paper's
matrix-semi-invariant presentation and classical ternary-cubic
invariant theorem are known inputs; an independent re-proof
of those earlier classification results is not claimed.

\begin{proof}[Proof of the positive Kronecker formula]
In characteristic zero, reorder the tensor factors in the
$N$-fold tensor power of $E^*\otimes F^*\otimes G^*$
and take diagonal $S_N$ invariants. Schur--Weyl decomposition
in the three factors gives
\[
 \operatorname{Sym}^N(E^*\otimes F^*\otimes G^*)
  =\bigoplus_{\alpha,\beta,\gamma\vdash N}
       (S_\alpha E^*\otimes S_\beta F^*\otimes S_\gamma G^*)
                      ^{\oplus g(\alpha,\beta,\gamma)},          \tag{A-026.4}
\]
where only partitions with at most three rows occur. Here
the coefficient is the invariant dimension of the three
Specht factors; their characteristic-zero self-duality gives
the usual Kronecker coefficient.

An irreducible polynomial $\mathrm{GL}_3$ module, or its
dual, restricts trivially to $\mathrm{SL}_3$ precisely when
its highest weight has all three parts equal. Such a module
is a one-dimensional determinant power and has degree $3h$.
Taking the three special-linear invariant spaces in
(A-026.4) therefore gives
\[
           \dim\mathcal A_{3h}
                      =g((h,h,h),(h,h,h),(h,h,h)),                \tag{A-026.5}
\]
and invariant spaces in all other degrees vanish.
The fixed choices of volume forms identify the determinant
lines with the scalar field; no dimension factors are lost.

By (A-026.3), the monomials
\[
                        I_6^a I_9^b I_{12}^c
\]
are a basis of $\mathcal A$. Their ordinary degree is
$6a+9b+12c=3(2a+3b+4c)$. Taking the degree-$3h$ basis
proves (A-026.1), including every multiplicity copy.
The positive indexing set is thus explicit, and multiplication
of its basis elements adds the three exponent vectors with
structure constant one. In particular this is more data
than an effective class in a representation $K_0$.
\end{proof}

\begin{proof}[Proof of the closed form]
Let $\epsilon\in\{0,1\}$ be the parity of $h$.
In (A-026.1), $b$ has that same parity. Write
$b=\epsilon+2j$ and put $M=(h-3\epsilon)/2$.
For $M\geq0$, the equation becomes
\[
                    a+2c+3j=M,
\]
so its number of solutions is
\[
           p(M)=\sum_{j=0}^{\lfloor M/3\rfloor}
                         \left(\left\lfloor\frac{M-3j}{2}\right\rfloor+1\right).
\]
Write $M=6q+s$, $0\leq s\leq5$. Pairing the terms
with indices $2i$ and $2i+1$ gives contribution $M-6i$.
For $s=0,1,2$ there are $q$ complete pairs and one
unpaired final term $\lfloor s/2\rfloor+1$; for
$s=3,4,5$ there are $q+1$ complete pairs. Consequently
\[
 p(M)=3q^2+(s+3)q+c_s,
           \qquad(c_0,c_1,c_2,c_3,c_4,c_5)=(1,1,2,3,4,5).
\]
This equals $\lfloor(M^2+6M+12)/12\rfloor$ in all
six residues. Substitution of $M=h/2$ for even $h$ and
$M=(h-3)/2$ for odd $h\geq3$ gives (A-026.2).
At $h=1$ the solution set is empty and the displayed odd
formula gives zero; at $h=0$ both formulas and the empty
Kronecker convention give one. This checks all endpoints.
The closed form requires only integer arithmetic on numbers
whose bit length is polynomial in the bit length of $h$.
\end{proof}

\paragraph{A primitive-lattice saturation failure inside an actual Kronecker ring.}
Grade $\mathcal A$ by $h=N/3$. Its three generator weights
are $2,3,4$. The support is exactly
\[
              S=\{0\}\cup\{2,3,4,\ldots\}.                     \tag{A-026.6}
\]
Indeed even $h\geq2$ uses $a=h/2,b=c=0$, and odd
$h\geq3$ uses $b=1,a=(h-3)/2,c=0$. No nonnegative
combination of $2,3,4$ is one. Since $3-2=1$, the
generated lattice is the full primitive lattice $\mathbb Z$.
Nevertheless $1\notin S$ while $2\cdot1\in S$.
Thus ordinary saturation fails even in this primitive lattice.
Equivalently
\[
                     g((1^3),(1^3),(1^3))=0,
             \qquad g((2^3),(2^3),(2^3))=1.
\]
The polynomial invariant algebra is a normal domain, and its
affine spectrum is smooth three-dimensional affine space.
These geometric properties do not repair the missing graded
piece. There is no conflict with A-019: its one-dimensional
graded-piece hypothesis fails, already at $h=4$, where
$I_6^2$ and $I_{12}$ are independent and the coefficient
is two. In fact at $h=4t$ the monomials with
$b=0,c=0,\ldots,t,a=2t-2c$ give at least $t+1$
independent vectors, so the multiplicities are unbounded.

In the finer exponent grading the support is the saturated
monoid $\mathbb Z_{\geq0}^3$. Its projection
$(a,b,c)\mapsto2a+3b+4c$ has the support (A-026.6).
This is an actual Kronecker realization of the failure of
saturation under a grading projection, not an unrelated
weighted-ring analogy.

\paragraph{What does and does not extend from A-024.}
For $r\geq3$, let
\[
       \mathcal R_r=\operatorname{Sym}
             (\mathbb C^3\otimes\mathbb C^3\otimes\mathbb C^r)
                                ^{\mathrm{SL}_3\times\mathrm{SL}_3}.
\]
Its $S_\nu\mathbb C^r$ multiplicity in degree $3h$ is
$g(\nu,(h^3),(h^3))$, by the same Cauchy argument.
Taking $h=4,\nu=(4,4,4)$ and using (A-026.1) gives
multiplicity two. Thus its highest-weight algebra has a
two-dimensional weight space and cannot satisfy the
multiplicity-free hypothesis of A-024. No classification
or nonpolynomiality claim about the full highest-weight
algebra is inferred. The equal-rectangle subalgebra in
(A-026.3) is precisely the positive specialization proved here.

\paragraph{Grothendieck dictionary and source limits.}
The three special-linear group schemes and their tensor
representation are defined over $\mathbb Z$; their invariant
spaces are kernels of comodule maps as in SGA 3, I,
Proposition 4.7.2, pp.~34--35. The determinant-line
condition is a Schur-functor calculation, with SGA 6, VI,
Theorem 3.3, pp.~392--393, and the exact Schur--Weyl
background in Macdonald, I, Appendix A(5.3)--(5.6),
pp.~155--156. The scalar and exponent gradings are the
torus gradings of SGA 3, I, Proposition 4.7.3, pp.~35--36.
The extra invariant-ring theorem (A-026.3) is essential;
it does not follow from these $K_0$ operations alone.

The primary source is available at
\url{https://arxiv.org/pdf/1101.3178}. Its characteristic-zero
statement is all that is used here. We make no assertion
that these invariant generators or their decomposition commute
with arbitrary integral base change. This positive family
and its saturation defect do not compare determinant and
permanent evaluation kernels and assert no additional
actual GCT multiplicity obstruction.
